% !TeX root = ../tcc.tex
% Detalhes de auditoria da busca e da triagem, retirados do Capítulo 3 na síntese do capítulo.
\chapter{Registro detalhado da busca e da triagem}
\label{apend:registro-busca}

Este apêndice reúne os registros da busca assistida e da triagem que
sustentam as quantidades apresentadas na Seção~\ref{sec:execucao-busca-selecao}:
o \textit{prompt} submetido às ferramentas, as consultas exibidas pelo
Consensus, os critérios avaliados pelo Jev Reviewer e a comparação entre
as sugestões automáticas e as decisões registradas no Evidentia.

\section{\textit{Prompt} submetido ao Consensus e ao Claude}

Como o Consensus recupera publicações por similaridade semântica, e não
apenas por correspondência booleana, a solicitação reuniu em inglês a
questão principal, as oito questões secundárias e os termos da
Tabela~\ref{tab:picoc-termos}. O Quadro~\ref{qua:prompt-consensus}
reproduz o texto submetido, que também foi usado na busca complementar
no Claude.

\lstinputlisting[
	caption={\textit{Prompt} estruturado submetido ao Consensus e ao Claude},
	label={qua:prompt-consensus},
	basicstyle=\scriptsize\ttfamily,
	breaklines=true,
	breakatwhitespace=true,
	columns=flexible,
	keepspaces=true,
	frame=single,
	rulecolor=\color{black!30},
	backgroundcolor=\color{black!2},
	xleftmargin=0.5em,
	xrightmargin=0.5em
]{editaveis/prompt-consensus.txt}
\vspace{-0.2cm}\noindent{\footnotesize\textbf{Fonte:} Autores.}\vspace{0.4cm}

\section{Consultas apresentadas pelo Consensus}

Na busca de 25 de setembro de 2026, feita no modo \textit{Deep} (19
etapas), o Consensus gerou 11 consultas em quatro rodadas
(Tabela~\ref{tab:consultas-consensus}). Os indicadores são saídas internas
da ferramenta e não correspondem a contagens auditadas de estudos únicos.

\begingroup
\small
\begin{longtable}{|p{0.08\textwidth}|p{0.68\textwidth}|r|}
	\caption{Consultas e indicadores apresentados pelo Consensus}\label{tab:consultas-consensus} \\ \hline
	Rodada          & Consulta                                                                                                                     & Indicador \\ \hline
	\endfirsthead
	\hline Rodada & Consulta                                                                                                                     & Indicador \\ \hline
	\endhead
	1               & deep learning convolutional neural network leprosy Hansen disease skin lesion classification detection diagnosis             & 95,3 mil  \\ \hline
	1               & U-Net ResNet EfficientNet Vision Transformer skin lesion segmentation leprosy hanseniasis dermatological pattern recognition & 2,3 mil   \\ \hline
	1               & machine learning model performance evaluation accuracy AUC F1 sensitivity specificity leprosy skin image diagnosis           & 63,5 mil  \\ \hline
	1               & skin tone fairness demographic subgroup performance dermatology AI model external validation generalization cross-dataset    & 67,4 mil  \\ \hline
	2               & leprosy skin image dataset Brazil public health AI diagnosis computer vision                                                 & 26,2 mil  \\ \hline
	2               & transfer learning self-supervised semi-supervised learning skin neglected tropical disease leprosy image classification      & 36,6 mil  \\ \hline
	2               & WHO skin neglected tropical diseases mobile application telemedicine leprosy detection low-resource settings                 & 63,4 mil  \\ \hline
	3               & explainability XAI Grad-CAM interpretable AI leprosy neglected tropical disease skin lesion diagnosis                        & 1,6 mil   \\ \hline
	3               & histopathological image analysis acid-fast bacilli detection deep learning leprosy nerve biopsy microscopy                   & 17,6 mil  \\ \hline
	4               & leprosy skin image dataset multimodal clinical metadata paucibacillary multibacillary classification AI model                & 3 mil     \\ \hline
	4               & reporting standards quality assessment QUADAS CONSORT-AI CLEAR dermatology AI model evaluation systematic review             & 47,8 mil  \\ \hline
\end{longtable}
\endgroup
\fonte{Autores.}

A plataforma descreve \textit{Retrieved} (425,4 mil) como o total de
trabalhos encontrados nas buscas, \textit{Eligible} (mil) como os
considerados relevantes após a deduplicação da ferramenta e
\textit{Included} (100) como os selecionados para a análise final, sem
esclarecer integralmente os critérios internos. Após a segunda rodada, o
histórico indica expansão por grafo de citações a partir de dez trabalhos
iniciais, com 662 trabalhos associados; após a terceira, a leitura de 12
resumos; antes da síntese, a leitura de resumos e PDFs de 100 trabalhos.
A busca também se estendeu a temas relacionados, como dermatologia geral,
equidade, saúde móvel, microscopia, histopatologia e diretrizes de relato.
A presença de uma referência na saída da ferramenta não significa que ela
atenda aos critérios da revisão: os 100 registros foram salvos pelos
autores e submetidos à deduplicação e à triagem.

\section{Critérios avaliados pelo Jev Reviewer e decisões de triagem}

O Jev Reviewer (\texttt{jev-1.13.0}) acompanha cada sugestão de valores
de \texttt{fails} associados a cinco critérios, que não substituem os
critérios de inclusão e exclusão da revisão
(Subseção~\ref{subsec:criterios-selecao}):
\begin{enumerate}
	\item abordar a hanseníase como condição investigada ou avaliada,
	      isoladamente ou entre outras doenças, e não apenas em menção
	      introdutória;
	\item utilizar fotografias clínicas ou imagens dermatoscópicas de
	      manifestações cutâneas humanas associadas à hanseníase, ou revisar
	      estudos que as utilizam;
	\item aplicar ou analisar métodos de visão computacional, aprendizado de
	      máquina ou aprendizado profundo para detectar, classificar,
	      identificar ou segmentar essas manifestações, incluindo revisões
	      dessas aplicações;
	\item relatar ou sintetizar evidências de avaliação desses métodos, como
	      experimentos, testes, validação, comparação de desempenho ou
	      resultados de estudos;
	\item não se limitar à inteligência artificial generativa como
	      intervenção.
\end{enumerate}

O cruzamento por título entre as sugestões do Jev Reviewer e as decisões
do Evidentia encontrou correspondência para os 119 registros. Entre os 27
incluídos no Evidentia, três tinham sugestão de inclusão, 11 sugestão
incerta e 13 sugestão de exclusão. Entre os 91 excluídos, 87 tinham
sugestão de exclusão e quatro sugestão incerta. O registro intermediário
tinha sugestão de exclusão. Os valores de \texttt{fails} não registram a
justificativa da decisão dos pesquisadores.

Os estados intermediários da triagem registram, em uma etapa, 13
inclusões, 32 exclusões, 15 decisões intermediárias e 59 registros sem
decisão; em outra, 20 inclusões, 40 exclusões e os mesmos 59 sem decisão,
com as 15 intermediárias convertidas em sete inclusões e oito exclusões.
As decisões preenchidas têm data de 28 de setembro de 2026, sem horário.
No resultado consolidado, 117 registros têm decisão consolidada
automática, segundo o campo correspondente do Evidentia; o registro~\estudo{9}
\cite{souza2021} recebeu duas exclusões individuais e foi consolidado
manualmente como intermediário; e o registro~\estudo{90} \cite{baweja2016}
recebeu uma inclusão e uma exclusão, conflito resolvido manualmente com
exclusão. Esse mecanismo de consolidação não identifica uma revisão
independente de cada artigo pelos dois pesquisadores.

\section{Observações sobre os critérios de seleção}

A formulação preservada do critério CI03 foi recebida truncada após
``resultados de estu''; a redação apresentada na
Subseção~\ref{subsec:criterios-selecao} resume a parte legível, e o texto
integral ainda deve ser confirmado. A verificação de disponibilidade e de
idioma do texto completo (CI01 e CE03) pode exigir consulta além do título
e do resumo, e a data e a duplicidade (CE01 e CE02) dependem também dos
metadados.
